% SPDX-License-Identifier: GPL-3.0-or-later
% Copyright (C) 2026 Matías Saavedra
%
% Design notes for Open Origamizer, one section per milestone.
% Build: latexmk -pdf -cd docs/notes.tex  (see CLAUDE.md, Commands)

\documentclass[11pt]{article}

\usepackage[T1]{fontenc}
\usepackage[utf8]{inputenc}
\usepackage{lmodern}
\usepackage{microtype}
\usepackage[margin=2.5cm]{geometry}
\usepackage{amsmath,amssymb,amsthm}
\usepackage{booktabs}
\usepackage{enumitem}
\usepackage[hidelinks]{hyperref}

% Deterministic output (see tvcg2010-notes.tex).
\pdfinfoomitdate=1
\pdftrailerid{}
\setlength{\emergencystretch}{2em}

\hypersetup{pdftitle={Open Origamizer: design notes}, pdfauthor={Open Origamizer project}}

\newcommand{\VTM}{\mathrm{VTM}}

\theoremstyle{definition}
\newtheorem{claim}{Claim}
\newtheorem*{remark}{Remark}

\title{Open Origamizer: design notes}
\author{Open Origamizer project}
\date{24 November 2026}

\begin{document}
\maketitle

\begin{abstract}
Running notes on design decisions that go beyond the papers, one section per
milestone. Notes on the papers themselves are in \texttt{tvcg2010-notes}. This
version covers M1 (preparing the input mesh and cutting it into a disk) and the decisions made so far for M2 (layout).
\end{abstract}

\section{M1: preparing and cutting the mesh}

Code: \texttt{include/origamizer/mesh.hpp} and \texttt{cut.hpp}; tests in
\texttt{tests/cpp} and \texttt{tests/python}.

\subsection{Conventions}

\begin{itemize}[nosep]
  \item Faces are listed counterclockwise around their normal, and the normal
    points to the \emph{clean side}; tucks go to the other side.
  \item A closed mesh is oriented outward, so its tucks go inside: if its
    signed volume is negative, all faces are reversed and the caller is told.
    A mesh with boundary keeps its orientation, which is the user's choice of
    clean side (TVCG10 Fig.~20).
  \item Preparing a mesh keeps the vertex numbering of the input. Face
    numbering changes only where nonconvex faces are split, and
    \texttt{face\_origin} maps every face back to its input face.
\end{itemize}

\subsection{Checks, and what is not repaired}

Before a mesh is built, the polygon soup must have no degenerate faces (fewer
than three vertices or a repeated vertex), no duplicate faces (the same vertex
set, in either orientation), no unused vertices, no distinct vertices at
exactly the same position, no edge shared by more than two faces, and no
vertex whose faces form more than one fan. It must be one component, joined
through edges, and its faces must be consistently oriented. Every violation is
reported, with indices; nothing is fixed silently. The one exception is
reorientation of an orientable soup, which happens only on request.

Each face must also be planar (no vertex farther than $10^{-6}$ face diameters
from the best-fit plane), have positive area (more than $10^{-12}$ times the
squared diameter), and be simple in its plane. Nonplanar faces are an error
rather than being triangulated: the choice of diagonal changes the surface, so
it belongs to the user.

The coincident-vertex check is disabled only when building a cut surface,
where the copies of a split vertex share its position.

\subsection{Convex facets}

The method needs convex facets: a concave corner violates the necessary
conditions T13 and T14 of TVCG10 (Sec.~5.2), and SoCG17 assumes convex facets.
TVCG10 triangulates every polygon, because extra edge-tucking molecules make
the conditions easier to meet. We split only nonconvex faces, into the fewest
convex pieces using diagonals only (CGAL's \texttt{optimal\_convex\_partition\_2},
Greene's dynamic program~\cite{greene83}), because every new diagonal becomes a
tuck across the face, and the fabrication track wants one rigid tile per facet.
Straight ($180^\circ$) corners are allowed. The pieces must cover the face
exactly; this is checked through their areas.

\subsection{Cutting into a disk}

The molecule mesh lives on a disk, so a surface that is closed, has handles, or
has several boundary loops is first cut along mesh edges (TVCG10 Sec.~2.2.1
and~8.1). After any cut the result is checked to be a disk: connected, one
boundary loop, Euler characteristic~1.

\paragraph{Algorithm.}
\begin{enumerate}[nosep]
  \item Flood the faces breadth first from a root face (TVCG10 Sec.~8.1). The
    edges the flooding crosses form a spanning tree $T^*$ of the dual graph and
    stay glued.
  \item Let $C$ be the interior edges not in $T^*$, and $G = C \cup \partial S$.
  \item While some vertex has exactly one edge of $G$, glue that edge back
    (remove it from $G$).
  \item If $S$ is closed and nothing is left, $S$ has genus~0: cut one shortest
    edge path between its two most curved vertices (largest angle defect, then
    largest distance, then smallest index), extended to two edges if it has
    only one.
  \item Split every vertex into one copy per fan of faces between cut edges, and
    rebuild the mesh.
\end{enumerate}

\begin{claim}
Cutting $S$ along all of $C$ gives a disk.
\end{claim}
\begin{proof}[Sketch]
In the cut surface, faces are glued only along the edges of $T^*$, a tree. Gluing
polygons one at a time along a single edge of the disk built so far keeps a
disk.
\end{proof}

\begin{claim}
Step~3 keeps a disk, and never removes a boundary edge.
\end{claim}
\begin{proof}[Sketch]
A vertex of the original boundary has two boundary edges in $G$, which are
never removed, so its degree stays at least~2. A vertex $v$ with one edge
$e = vw$ of $G$ is therefore interior, and in the cut surface it is the tip of
a slit: the two copies of $e$ are consecutive boundary edges meeting at $v$.
Gluing them closes the slit from its tip, which keeps a disk as long as the
boundary has more than these two edges. The only case where it does not is a
closed surface cut along a single edge; that case is caught by step~4 and by
the check below.
\end{proof}

\paragraph{What remains.} If every vertex of $S$ touches $G$, the cyclomatic
number of $G$ is $|E(G)| - |V(G)| + 1 = (|C| + E_\partial) - V + 1$. With
$|C| = E - E_\partial - (F - 1)$ this is $2 - \chi = 2g + b$ for a connected
orientable surface of genus~$g$ with $b$ boundary loops, and pruning does not
change it. So the pruned graph consists of the $b$ boundary loops, $2g$ loops
around the handles, and paths joining them. This corrects the earlier plan of
``one cut path per handle'': a torus needs two loops, not one.

\paragraph{One-edge slits.} If a cut edge has neither end split, by another cut
edge or by the boundary, its two copies join the same two vertices, so
rebuilding the mesh glues them back. This is why step~4 extends a one-edge path
(every two vertices of a tetrahedron are adjacent), and why
\texttt{cut\_along} rejects such an edge with an explicit message.

\paragraph{Splitting vertices.} A \emph{corner} is a halfedge $h$ with a face:
the corner of that face at the target of~$h$. For every glued interior edge,
the corners on its two sides at each of its endpoints are joined in a
union-find structure. Each class is one copy of a vertex, and copies of the
same vertex are numbered consecutively.

\subsection{Results on the examples}

\begin{table}[ht]
\centering
\small
\begin{tabular}{@{}lrrrrrrr@{}}
\toprule
Mesh & $V$ & $E$ & $F$ & $\chi$ & $b$ & $g$ & cut edges \\
\midrule
tetrahedron & 4 & 6 & 4 & 2 & 0 & 0 & 2 \\
cube & 8 & 12 & 6 & 2 & 0 & 0 & 3 \\
octahedron & 6 & 12 & 8 & 2 & 0 & 0 & 2 \\
L prism (caps split) & 12 & 20 & 10 & 2 & 0 & 0 & 3 \\
torus & 96 & 192 & 96 & 0 & 0 & 1 & 20 \\
open cylinder & 24 & 36 & 12 & 0 & 2 & 0 & 1 \\
saddle & 49 & 120 & 72 & 1 & 1 & 0 & 0 \\
Stanford bunny (372 triangles) & 195 & 570 & 372 & $-3$ & 5 & 0 & 45 \\
\bottomrule
\end{tabular}
\caption{Topology and automatic cuts. The torus is cut along one loop around the
hole (12 edges) and one around the tube (8 edges). The bunny's five boundary
loops are the holes of the scan.}
\end{table}

\begin{remark}[Open for M2]
The cut is only guaranteed to be valid; its quality matters a lot. TVCG10
reports a bunny efficiency of 0.007 with its automatic cut and 0.172 after
moving the boundary behind the ears, and notes that a boundary through sharp
points helps (Sec.~10); it also smooths ``jags'' out of the cut (Sec.~8.1).
Our paths between boundary loops follow the flooding tree rather than
shortest paths, which is why the bunny needs 45 cut edges, and the root face
is fixed at~0. Once M2 can measure a layout, compare cut strategies on these
meshes; for scanned meshes, also consider filling small holes.
\end{remark}

\subsection{Test data}

\texttt{examples/meshes/generate.py} writes the seven example meshes and checks
their orientation. \texttt{fetch\_bunny.py} downloads the Stanford bunny,
drops the vertices that no face uses, and simplifies it to at most 374
triangles with CGAL edge collapses (the size used in SoCG17 Fig.~1). The coarser
reconstructions in the Stanford archive are not 2-manifolds (duplicate faces,
edges shared by three faces), so the full one is simplified instead. The bunny
is not committed: the Stanford terms allow research use and free
redistribution but not commercial use without permission, which does not fit a
GPL repository.

\section{M2: layout}

Code so far: \texttt{molecules.hpp}, \texttt{layout.hpp} (rotation step,
translation step, reconstruction), and the private solver wrapper
\texttt{quadratic\_program}. This section records the decisions as they are
made; equation tags T1 to T19 refer to \texttt{tvcg2010-notes}.

\subsection{Molecule mesh}

On the cut disk, every interior edge gets an edge-tucking molecule, oriented
from its smaller vertex index to the larger, and every boundary vertex gets a
boundary molecule. The \emph{star} of a vertex lists its molecules
counterclockwise, seen from the clean side; slot $n$ carries the sign of
$\theta(v, j_n)$ relative to the molecule's orientation and the sector
$\alpha(v, j_n)$, the corner of the polygon between slots $n-1$ and~$n$. With
counterclockwise faces, the corner of the face of an incoming halfedge $h$
lies between \texttt{edge(next($h$))} (clockwise side) and
\texttt{edge($h$)} (counterclockwise side), and the next face counterclockwise
is that of \texttt{prev(opposite($h$))}; CGAL's
\texttt{halfedges\_around\_target} runs clockwise, so the code walks the stars
itself. A test on a flat hexagonal fan checks the order. Built in, and checked
on construction: each edge molecule appears in the stars of both its ends with
opposite signs, and each boundary vertex has one boundary molecule. Tested:
the defects of a disk sum to $2\pi + \pi N_\partial$, with $N_\partial$ the
number of boundary vertices (the Gauss-Bonnet identity of
\texttt{tvcg2010-notes}, Sec.~4).

\subsection{Boundary digons (default changed)}

T8 makes every vertex-tucking molecule a convex polygon: each exterior angle
$\Theta_m$ lies in $[0, \pi)$, and by T1 (with T3) they sum to $2\pi$. So a star
needs at least three slots. In the plain reading of the boundary molecules, a
boundary vertex with $K$ faces has $K$ slots ($K-1$ edge molecules and the
boundary molecule), so any boundary vertex with one or two faces makes the
rotation step infeasible. Every example has such vertices:

\begin{center}
\small
\begin{tabular}{@{}lrr@{}}
\toprule
Disk & boundary vertices & with one or two faces \\
\midrule
tetrahedron & 4 & 2 \\
cube & 6 & 4 \\
octahedron & 4 & 2 \\
L prism & 6 & 3 \\
torus & 40 & 40 \\
open cylinder & 26 & 26 \\
saddle & 24 & 4 \\
Stanford bunny & 114 & 44 \\
\bottomrule
\end{tabular}
\end{center}

This explains TVCG10 Sec.~5.1, which adds, along every boundary edge, a
zero-area digon glued to the face by an extra edge-tucking molecule ``to avoid
overconstraints by convex paper conditions''. Every boundary star gains two
slots: counterclockwise, the edge molecules between the faces (sectors
$c_1, \dots, c_{K-1}$), then the digon molecule of the counterclockwise boundary
edge (sector $c_K$, the last face's corner), the boundary molecule (sector~0,
a digon corner), and the digon molecule of the clockwise boundary edge
(sector~0). The digon molecules are shared by the stars of both ends of their
edge. The digons are therefore on by default (they were first proposed as an
option); \texttt{MoleculeOptions::boundary\_digons} turns them off. This
settles Open question~1 of \texttt{tvcg2010-notes} for the equations; how
digons appear in the 2D layout comes with the reconstruction.

\subsection{Rotation step: results}

With the digons, the rotation program (below) is feasible on all seven
committed examples, in 0.6\,s for all seven together. On every solution the
residuals of T15, T4, T6, and T8 are below $10^{-6}$, and the boundary
molecules turn by exactly one full turn,
$\sum_{\partial}(\theta(i,o) - \pi) = 2\pi$, as the Gauss-Bonnet argument of
\texttt{tvcg2010-notes} predicts. The rotation equations have full row rank
($N_{\mathrm{vert}}$) on a disk: their matrix is the signed incidence matrix of
a connected graph (rank $N_{\mathrm{vert}} - 1$) plus the boundary molecules'
columns. This answers the rotation half of Open question~5.

\subsection{Solver}

Both mapping steps are linear constraint problems apart from T19, so they are
solved as linear or quadratic programs rather than by the paper's projected
gradient descent (TVCG10 Sec.~6.2 and~6.3). A solver gives exact answers to
its tolerance and can prove that no solution exists, which tells the pipeline
when to split edges or remap; a penalty method only shows a penalty that stops
decreasing.

The user first preferred CGAL's exact LP/QP solver, to avoid a new dependency.
On synthetic problems shaped like the translation step (sparse equality rows,
two per vertex; every width in two vertices; widths at least 0.01) it was far
too slow:

\begin{center}
\small
\begin{tabular}{@{}lrrrr@{}}
\toprule
Problem & variables & rows & CGAL (exact) & HiGHS \\
\midrule
cube-sized & 30 & 24 & 0.05\,s & 0.009\,s \\
torus-sized & 300 & 220 & over 10\,min (stopped) & 0.009\,s \\
bunny-sized & 600 & 480 & not run & 0.02\,s \\
\bottomrule
\end{tabular}
\end{center}

CGAL's solver needs exact arithmetic (with floating point it aborts on its own
basis-inverse check), and its documentation targets problems with few
variables or few constraints. We therefore use HiGHS~\cite{highs}, MIT
licensed, pinned at v1.15.1 and built from source by CMake (about 100\,s on a
clean build). It is called through a small private wrapper, so its types do
not appear in public headers.

\subsection{Solver accuracy: polishing and a KKT check}

HiGHS sends quadratic programs to its active-set solver, which adds $10^{-7}$
to the diagonal of the Hessian (option \texttt{qp\_regularization\_value})
and stops at hard-coded thresholds, so its optima are accurate to about
$10^{-8}$. That noise matters here because the exact optima have structure.
On the cube the exact rotation angles are multiples of $\pi/10$, and HiGHS
returned, for example, $-0.099999999\,\pi$ and $1.5\times10^{-8}\,\pi$ for
$-0.1\,\pi$ and~$0$. In the translation rows such angles give coefficients of
$2\times10^{-8}$ where the exact value is~0, and exact dependencies between
rows become near-dependencies (the smallest singular value of the cylinder's
$C_w$ was $2.7\times10^{-9}$). This did real damage downstream: objective~A
called the cube feasible, with widths near $10^8$ along a direction that only
the noise opened; HiGHS's simplex stalled on the cube's phase-1 program with a
dual infeasibility of $2.8\times10^{-7}$; and its QP solver returned, on the
cylinder's phase~2, a point it called optimal that fails the optimality
conditions. Two measures in the wrapper now cover every program:
\begin{itemize}[nosep]
  \item \emph{Polishing.} HiGHS's basis status gives the active set: the
    variables at a bound and the rows at a bound (equality rows always count).
    The wrapper fixes those variables, keeps those rows as equations, and
    solves the KKT system of what remains,
    $D_{FF}x_F - A_{RF}^\top y_R = -c_F$ and $A_{RF}x_F = b_R - A_{RX}x_X$,
    with a sparse QR (Eigen), accepting the solve only if its residual is
    below $10^{-9}$. This is the polishing step of solvers such as OSQP. The
    cube's angles then come out as multiples of $\pi/10$ to within
    $10^{-16}$, and the cylinder's as exactly $0$, $\pi$, and $3\pi/2$.
  \item \emph{KKT check.} A point counts as optimal only if it satisfies the
    KKT conditions, checked independently of HiGHS: primal feasibility, and
    with $z = Dx + c - A^\top y$, every multiplier of the right sign at its
    bound and zero away from it (relative tolerance $10^{-6}$). The polished
    point is checked first, then HiGHS's own; if neither passes, the solve
    fails with a message instead of returning a wrong optimum.
\end{itemize}
Two smaller rules: the wrapper drops matrix entries of magnitude at most
$10^{-9}$ (HiGHS's \texttt{small\_matrix\_value}; HiGHS drops them too, but
reports a warning, which the wrapper treats as a rejection), and the
translation step works in units of $\ell_{\mathrm{ave}}$, so that the solver's
absolute tolerances ($10^{-7}$) stay small next to the widths on meshes of any
size.

\subsection{Objectives}

The constraints of each step leave freedom, so each step also optimizes
something. Strict inequalities become margins: T6 becomes
$|\theta(i,j)| \le \pi - \epsilon_\theta$ and T8 becomes
$0 \le \Theta_m \le \pi - \epsilon_\Theta$, with
$\epsilon_\theta = \epsilon_\Theta = 10^{-3}$ by default.

\paragraph{Rotation step.} Minimize $\sum \theta_n^2$ over the edge and digon
molecules, subject to T15 ($C_\theta\boldsymbol\theta = \mathbf g$), T4, T6,
and T8. The paper starts its descent from the minimum-norm solution
$C_\theta^+\mathbf g$; this picks the feasible angles closest to rectangular
tucks. Boundary molecules are left out of the sum, since T4 keeps them at or
above~$\pi$. If the program is infeasible, an elastic linear program keeps
T15 and minimizes the total violation of the inequalities: every bound of T4
and T6 and every T8 row gets a slack $s \ge 0$ that widens it on both sides,
and $\sum s$ is minimized. It returns the least-violation angles and the
broken inequalities, largest first. (The notes first planned squared
violations; the sum is a linear program, the same choice as phase~1 of the
translation step below, and it is zero exactly when the squares are.)

\paragraph{Translation step, objective A (minimum width).} Minimize
$\sum_n w_n$ subject to T17 ($C_w\mathbf w = \mathbf b$) and
$w_n \ge w_{\min}$ for every molecule (T5 and T7 with a margin), a linear
program. $w_{\min}$ defaults to $0.01\,\ell_{\mathrm{ave}}$, a hundredth of the
mean edge length; for fabrication it is the smallest gap that can be made,
which is OriStitch's minimum separation as a hard bound.

\paragraph{Translation step, objective B (the paper's penalty).} TVCG10
Sec.~6.3 minimizes $E_w = c_0E_{(5,7)} + c_1E_{(19)} + c_2|\Delta\mathbf w|^2$.
We solve it in two phases. Phase~1 minimizes the total negative width
$\sum_n s_n$, with slack variables $w_n + s_n \ge 0$ and $s_n \ge 0$, subject
to T17: a linear program whose optimum is zero exactly when T5 and T7 can
hold, which is also exactly when the paper's $E_{(5,7)}$, the sum of the
squared violations, can vanish. (Changed on 24 November 2026: phase~1 first
minimized $\sum_n s_n^2$, as accepted. In that program the widths are free
variables with zero curvature, and HiGHS's QP solver failed on it: an
``optimum'' with a violation of 3.84 on the cube where the same angles
admitted nonnegative widths, ``solve error'' on three examples, and an
internal assertion on the tetrahedron. The change does not alter which
layouts are feasible; for an infeasible one, phase~1 now returns the widths
of least total violation instead of least squared violation.) If it is zero,
phase~2 finds
the widths closest to a reference, $\min |\mathbf w - \mathbf w_{\mathrm{ref}}|^2$
subject to T17 and $\mathbf w \ge 0$, where $\mathbf w_{\mathrm{ref}}$ is the
paper's starting point $C_w^+\mathbf b$ by default, or widths requested by the
user (the role of the $c_2$ term). Otherwise phase~1's solution is returned
with its violations reported, as the paper's penalty would leave them. The
$c_1E_{(19)}$ term (weight $\ell_{\mathrm{ave}}^2$) joins once T19 is
implemented.

\subsection{Translation step: rank and results}

Before solving, the translation step computes the minimum-norm least-squares
solution of T17 (the paper's starting point $C_w^+\mathbf b$, with Eigen's
complete orthogonal decomposition) and stops with a message if it leaves a
residual: then T17 has no solution, nonnegative or not. Otherwise the programs
get a maximal independent subset of the rows (column-pivoted QR of $C_w^\top$,
relative threshold $10^{-10}$), because an active-set solver needs independent
equality rows; the final check of the widths covers every row.

The rows are dependent on four of the seven examples. The paper counts
$N_e - 2N_v$ degrees of freedom for the widths (TVCG10 Sec.~6.3), which
assumes full row rank:

\begin{center}
\small
\begin{tabular}{@{}lrrrcc@{}}
\toprule
Disk & $N_v$ & molecules & rank of $C_w$ & objective A & objective B, phase~1 \\
\midrule
tetrahedron & 5 & 12 & 9 & feasible & 0 \\
cube & 10 & 21 & 20 & infeasible & $11.46\,\ell_{\mathrm{ave}}$ (9 negative) \\
octahedron & 7 & 18 & 13 & feasible & 0 \\
saddle & 49 & 144 & 98 & feasible & 0 \\
torus & 117 & 252 & 233 & feasible & 0 \\
open cylinder & 26 & 63 & 51 & feasible & 0 \\
L prism & 14 & 29 & 28 & infeasible & $8.16\,\ell_{\mathrm{ave}}$ (8 negative) \\
\bottomrule
\end{tabular}
\end{center}

\paragraph{Where the dependency comes from (partial).} Rotate the equations of
each vertex into the plane of the layout. Around $\VTM(v)$, side $n$ points
along a unit vector $d_{v,n}$. The two bases of an edge-tucking molecule
between $u$ and $v$ are parallel, and $\VTM(u)$ and $\VTM(v)$ traverse them in
opposite directions, so the molecule's width enters the closure of $v$ along
some $d$ and the closure of $u$ along $-d$. A left null vector therefore
assigns a vector $y_v$ to each vertex with $(y_v - y_u)\perp d$ for every edge
or digon molecule between $u$ and $v$, and $y_v \perp d_{v,o}$ for every
boundary molecule. For a triangulated disk these
are $3N_v - 3$ conditions on $2N_v$ unknowns, so dependencies are not generic:
they need special angles, and the four examples with one are highly
symmetric. On all four, $\mathbf b$ lies in the range of $C_w$ (the step
succeeds), so the equations are redundant there, not contradictory. Whether
that holds in general is open; the consistency check above would report it.

\paragraph{Infeasible translations.} On the cube and the L prism no
nonnegative widths exist for the angles of the rotation step, and the KKT
check certifies phase~1's optimum (for a linear program, the KKT conditions
are sufficient). This is the failure of TVCG10 Sec.~6.4: fixing
$\boldsymbol\theta$ first discards freedom. Remapping (below) repairs the
cube; the L prism needs a longer cut. Both steps take milliseconds: on the
torus, 21\,ms for the rotation step, 21\,ms for objective A, and 45\,ms for
objective B (Release build).

\subsection{Remapping}

TVCG10 Sec.~6.5 changes the angles when the translation step fails: draw the
molecule graph without crossings by Tutte's barycentric method with a convex
boundary, measure the angles of that drawing, fit $\boldsymbol\theta$ to them
through T3 by least squares, and restart the rotation step from there. Our
reading of the graph: its nodes are the polygons (faces and digons), its edges
the molecules, and its faces the vertex-tucking molecules. That is the
homogeneous drawing of TVCG10 Sec.~6.4 with the polygons shrunk to points, so
a Tutte drawing, whose inner faces are convex, has exterior angles that
satisfy T8 and close up, and its edge lengths are positive widths for those
angles. The boundary molecules join consecutive digons and form the outer
cycle. \texttt{remap\_angles} places the digons evenly on a circle, solves
the barycentric system with uniform weights (sparse Cholesky), flips the sign
of every angle if the faces come out clockwise, and fits the angles with a
sparse QR; it needs the boundary digons.

The paper's rotation step is a descent from the current angles and stops at
the first feasible point, so a restart stays near the drawing. Ours returns
the feasible angles closest to a reference, $\sum(\theta_n - \theta_n^{\mathrm{ref}})^2$
over the edge and digon molecules, with $\boldsymbol\theta^{\mathrm{ref}} = 0$ by
default and the fitted angles for remapping. The drawing does not depend on
the angles, so repeating the step gives the same reference; we remap once.
On the cube the fitted angles miss T15 by 0.18 and the rotation step moves
them by at most 0.018: the result admits nonnegative widths for every one of
the six root faces of the automatic cut, against none of the six without
remapping. The five examples that were already feasible stay feasible. On the
L prism the fit misses T15 by 0.34 and the inequalities by 0.62, and no root
face works, with or without remapping.

\subsection{Cut strategy}

M1's automatic cut prunes the cut graph to its essential part, so a closed
genus-0 surface gets a single short seed path (3 edges on the cube and the L
prism). TVCG10 Sec.~8.1 instead keeps the cut tree of its flooding,
simplified by removing ``branches and jags'', and Sec.~10 notes that adding
boundary edges adds freedom. \texttt{CutStrategy::kTree} keeps the whole cut
tree, every interior edge the flooding does not cross, which puts every
vertex on the cut. Results with objective~B (``ok'' means nonnegative widths
exist; numbers are phase~1's total negative width in mean edge lengths):

\begin{center}
\small
\begin{tabular}{@{}llrrll@{}}
\toprule
Mesh & cut & edges & boundary vertices & plain & remapped \\
\midrule
cube & minimal & 3 & 6 & 11.46 & ok \\
cube & tree & 7 & 14 & ok & ok \\
L prism & minimal & 3 & 6 & 8.16 & 8.17 \\
L prism & minimal, plus paths to its 2 saddles & 4 & 8 & 8.17 & 8.25 \\
L prism & tree & 11 & 22 & ok & ok \\
\bottomrule
\end{tabular}
\end{center}

The tree cut worked from each of the first three root faces tried. It costs
seams: every cut edge becomes two stretches of paper boundary that meet
again in the folded model, so the cube's 3 seams become 7. The default cut
stays \texttt{kMinimal}. Between the two lie partial trees (the paper's
simplification, or pruning the tree while the mapping stays feasible), to
compare once M4 measures layouts.

\subsection{Mapping driver}

\texttt{map\_surface} runs the steps in the order the results above suggest.
For each cut in the list \texttt{MappingOptions::cuts} (minimal, then tree),
it builds the molecule mesh, runs the rotation and translation steps, and if
the translation fails, remaps and runs both again; it stops at the first
success. It returns the
cut, the molecule mesh, the angles, the widths, and every attempt with the
message of the step that failed. With objective~B a failed mapping keeps
phase~1's widths, so the crossed molecules can be drawn. On the committed
examples the first attempt succeeds, except on the cube (minimal cut,
remapped: second attempt) and the L prism (tree cut without remapping: third
attempt). \texttt{examples/plot\_layout.py} draws the result for every
example.

\subsection{Overlap check}

TVCG10 Sec.~4.2.1 argues that with a convex outline the molecule mesh cannot
overlap unless a cell intersects itself or flips, and takes T4 to T8 as
sufficient conditions. \texttt{check\_overlap} verifies the conclusion on the
placed polygons rather than trusting the conditions: every edge- and
vertex-tucking molecule must be convex (no negative turning angle, skipping
sides of zero length) with nonnegative signed area, and the outline must turn
by exactly $2\pi$ with no reflex corner, so it is a simple convex polygon. As
a consistency check, the areas of the faces and molecules must add up to the
outline's area (Green's theorem makes this hold for any closed molecule mesh,
so it tests the outline and the cells, not the overlap). The outline comes
from the boundary stars: for each boundary vertex in order, its copies in the
digons before and after its boundary molecule, joined alternately by the
boundary molecule's side and by a digon. \texttt{map\_surface} reconstructs
every feasible layout and requires the check to pass; it does on all seven
examples, with every molecule of positive area under objective~A, while the
cube with the rotation step's angles and phase~1's widths fails it with
flipped molecules. Objective~B allows zero widths, and then molecules
degenerate to points or segments (on the open cylinder every width is zero:
the strip needs no tucks). A degenerate cell covers nothing, so the check
accepts cells whose area is within $10^{-9}\ell_{\mathrm{ave}}^2$ of zero and
tests convexity only on the others; the first version tested every cell and
misread the $\pm\pi$ turns of a segment as reflex corners.

\subsection{Crease conditions T9 and T19: measurements}

\texttt{crease\_conditions} evaluates T9 and T19 on the placed layout: for
every side of positive length of every vertex-tucking molecule, it finds the
Delaunay triangle of the molecule's corners on that side (CGAL) and compares
its apex angle $\varphi$ with $\gamma + \pi/2$ (T9) and its base angles with
$-\frac12\theta(v,j)$ (T19). On a boundary molecule's side there is no
neighboring molecule; we require $\varphi \le \pi/2$ there, so that the
Voronoi vertex stays on the paper. That is our extension, not a condition of
TVCG10. Counts with objective~A, sides of boundary molecules in parentheses:

\begin{center}
\small
\begin{tabular}{@{}lrrr@{}}
\toprule
Example & sides & T9 fails & T19 fails \\
\midrule
tetrahedron & 20 (4) & 0 (4) & 1 \\
cube & 36 (6) & 0 (6) & 4 \\
octahedron & 32 (4) & 0 (4) & 10 \\
saddle & 264 (24) & 0 (0) & 0 \\
torus & 464 (40) & 0 (6) & 32 \\
open cylinder & 100 (26) & 0 (0) & 0 \\
L prism & 84 (22) & 0 (1) & 4 \\
\bottomrule
\end{tabular}
\end{center}

T9 holds on every side shared with an edge-tucking molecule (objective~B
breaks it on two sides of the torus). On boundary sides it often fails: the
Delaunay triangle on a boundary molecule's side is obtuse, so the Voronoi
folding of TVCG10 Sec.~3.2 would put a crease vertex outside the paper. M3
has to fold boundary vertex-tucking molecules differently, or the layout
must keep those triangles acute; TVCG10 does not discuss it. T19 holds on
the saddle and the cylinder and fails on some sides of every convex example.
It is the condition under which edge splitting can satisfy T9 to T11
(TVCG10 Sec.~7.2), and the paper keeps it with the term $c_1E_{(19)}$ of the
translation step. Since the Delaunay triangles depend on the widths, that
term needs an iterative scheme (sequential linearization of the base angles
around the current widths is the natural candidate); it is the next step of
M2, followed by edge splitting.

\subsection{The Stanford bunny: the rotation step fails}

The bunny generated by \texttt{fetch\_bunny.py} (372 triangles from CGAL edge
collapses; not committed) makes the rotation step infeasible: the least total
violation is 0.408\,rad, all of it in T8 at two interior vertices (0.208 and
0.200\,rad). Both are strong saddles, with angle defects of $-1.80$ and
$-1.10$\,rad, next to sliver corners of $9^\circ$ to $11^\circ$; at the second,
two molecules already sit at the T6 bound. T8 is TVCG10's sufficient
condition for convex vertex-tucking molecules, which its Voronoi crease
patterns need, so the paper's rotation step would also stop with a positive
penalty on this mesh. Extending the cut through the failing vertices (shortest
paths to the existing cut, the paper's tool of Sec.~8.1) does not help: the
least violation stays at 0.408\,rad and moves to the copies of the same
vertices. Neither does the whole cut tree (175 cut edges, all 374 disk
vertices on the boundary): the rotation step stays infeasible. The mesh is
the problem: its smallest corner is $2.7^\circ$, 54
corners are below $15^\circ$, and one vertex has a defect of $-3.09$\,rad,
while no committed example has a corner below $32^\circ$. Edge collapses
optimize shape error, not triangle quality. A better test mesh (isotropic
remeshing to the same size) is the likely remedy; to be decided.

\subsection{Layout reconstruction}

\texttt{reconstruct\_layout} places the polygons from $\boldsymbol\theta$ and
$\mathbf w$. Every star slot records the polygon at its sector (faces first,
then digons), and every edge or digon molecule joins the polygon before it,
$P_0$, and the one after it, $P_1$, counterclockwise around its vertex~$i$.
With $u_0$ the unit direction of $P_0$'s copy of the leg $i \to j$, the
molecule's geometry (\texttt{tvcg2010-notes}, Fig.~1) gives $P_1$'s copy:
\[
  A_1 = A_0 + w(i,j)\,R\bigl(\tfrac12\theta + \tfrac12\pi\bigr)u_0, \qquad
  u_1 = R(\theta)\,u_0, \qquad w(i,j) = w_n + \delta(i,j),
\]
with $R$ a counterclockwise rotation. The first face is placed in its own
coordinates and the others breadth first, digons as segments. Then every side
of every vertex-tucking molecule is compared with T2 and T3: side $n$ of
$\VTM(v)$ must have length $w(v,j_n)$ and direction $\varphi_v + \psi_{v,n}$,
with $\varphi_v$ read off one molecule of the star. Tree molecules pass by
construction, so the comparison tests the other molecules and the boundary
molecules, independently of how \texttt{translation\_equations} assembles
T17. On all seven examples the largest error is below
$2\times10^{-13}\,\ell_{\mathrm{ave}}$, for the cube and the L prism with
phase~1's negative widths too, and a perturbed width is rejected. This
confirms the conventions of T2, T3, and N6, and that the infeasibility above
is real.

\subsection{Defaults accepted by the user}

\begin{itemize}[nosep]
  \item Boundary molecules follow our reading in \texttt{tvcg2010-notes}
    (Sec.~4): one per boundary vertex, standing for the outside of the paper,
    with the two boundary edges as its legs. The extra molecules that TVCG10
    Sec.~5.1 adds along the boundary are an option, first tested on a single
    triangle and a $2\times2$ grid. (Later changed: the digons are on by
    default; see above.)
  \item Convex faces stay whole. If the cube's quadrilaterals cannot be laid
    out, triangulating faces, as TVCG10 does by default, is the fallback.
  \item 2D conditions come first (T1 to T9, T19, edge splitting); the tuck
    proxy and the 3D conditions come last.
\end{itemize}

\begin{thebibliography}{9}
\bibitem{highs}
Q.~Huangfu and J.~A.~J. Hall. Parallelizing the dual revised simplex method.
\emph{Mathematical Programming Computation} 10:119-142, 2018. Software:
HiGHS, \url{https://highs.dev}.
\bibitem{tvcg10}
T.~Tachi. Origamizing polyhedral surfaces. \emph{IEEE Transactions on
Visualization and Computer Graphics} 16(2):298-311, 2010.

\bibitem{socg17}
E.~D. Demaine and T.~Tachi. Origamizer: A practical algorithm for folding any
polyhedron. \emph{SoCG 2017}, LIPIcs, pp.~34:1-34:15.

\bibitem{greene83}
D.~H. Greene. The decomposition of polygons into convex parts. In
F.~P. Preparata, editor, \emph{Computational Geometry}, Advances in Computing
Research~1, pp.~235-259. JAI Press, 1983.
\end{thebibliography}

\end{document}
